\begin{definition}
Сводимость по Куку  $R_1$ к $R_2$ - это некоторый полиномиальный алгоритм, который решает задачу $R_1$ при условии, что функция, которая находит ответ к задаче $R_2$, дана в качестве оракула, т.е. решение задачи $R_2$ можно получить за 1 шаг.
\end{definition}

\textbf{Задача поиска}\\
Дана $V(x,y)$. Нужно понять, существует ли $y$, такое что $V(x,y) = 1$.
\begin{theorem}
Пусть $A$ --- \textbf{NPC} задача. Пусть машина может распознавать, лежит ли слово $x$ в $A$ за 1 шаг. Тогда существует полиномиальный алгоритм, решающий задачу поиска для $A$ (то есть задача поиска сертификата сводится по Куку к задаче распознавания, есть ли он).
\end{theorem}
\begin{example}
Проверим для задачи поиска клики размера $k$. Возьмем произвольную вершину $v$ и удалим ее из графа. Если в данном графе сохранилась клика размера $k$, то продолжим искать ее в уменьшенном графе. Иначе -- вершина $v$ входит в эту клику, то есть можно исключить все вершины, не связанные с $v$. В сокращенном графе будем искать клику размера $k-1$. И вместе с $v$ она даст нужную клику.
\end{example}
